% FAC-EN-U0057
% Source: fac_body.tex lines 2043--2076.
% Authority: PDF physical pp. 56--57; printed pp. 251--252 (no. 59 complete).
% English source-first translation; formulas and mathematical structure source-aligned.

\phantomsection\label{fac:III:2:no:59}
\textbf{59. Graded $S$-module associated with an algebraic sheaf.}

Let $\mathcal{F}$ be an algebraic sheaf on $X$, and set

\phantomsection\label{fac:III:2:no:59:formula:graded-sections}
\[
\Gamma(\mathcal{F})=
\sum_{n\in\mathbf{Z}}\Gamma(\mathcal{F})_n,
\qquad
\Gamma(\mathcal{F})_n=\Gamma(X,\mathcal{F}(n)).
\]
The group $\Gamma(\mathcal{F})$ is graded; we shall endow it with the
structure of an $S$-module. Let $s\in\Gamma(X,\mathcal{F}(q))$ and
$P\in S_p$. We may identify $P$ with a section of $\mathcal{O}(p)$
(compare \hyperref[fac:III:2:no:54]{no.~54}); hence $P\otimes s$ is a
section of
\phantomsection\label{fac:III:2:no:59:formula:S-action}
\[
\mathcal{O}(p)\otimes\mathcal{F}(q)
=\mathcal{F}(q)(p)=\mathcal{F}(p+q),
\]
where we have used the isomorphisms of
\hyperref[fac:III:2:no:54]{no.~54}. We have thus defined a
section of $\mathcal{F}(p+q)$, which we denote by $P.s$ rather than
$P\otimes s$. The map $(P,s)\mapsto P.s$ gives
$\Gamma(\mathcal{F})$ the structure of an $S$-module compatible with
its grading.

One can also define $P.s$ in terms of its components over the $U_i$.
If the components of $s$ are
$s_i\in\Gamma(U_i,\mathcal{F})$, with
$s_i=(t_j^q/t_i^q).s_j$ on $U_i\cap U_j$, then
\phantomsection\label{fac:III:2:no:59:formula:component-action}
\[
(P.s)_i=(P/t_i^p).s_i.
\]
This is well defined, since $P/t_i^p$ is a regular function on $U_i$.

To compare the functors $\mathcal{A}(M)$ and $\Gamma(\mathcal{F})$, we
shall define two canonical homomorphisms:
\phantomsection\label{fac:III:2:no:59:formula:alpha-beta}
\[
\alpha:M\longrightarrow\Gamma(\mathcal{A}(M)),
\qquad
\beta:\mathcal{A}(\Gamma(\mathcal{F}))\longrightarrow\mathcal{F}.
\]

\phantomsection\label{fac:III:2:no:59:def:alpha}
\textsc{Definition of $\alpha$.}
Let $M$ be a graded $S$-module, and let $m\in M_0$ be a homogeneous
element of degree zero in $M$. The element $m/1$ is a well-defined
element of $M_x$, and varies continuously with $x\in X$; hence $m$
defines a section $\alpha(m)$ of $\mathcal{A}(M)$. If now $m$ is
homogeneous of degree $n$, then $m$ is homogeneous of degree zero in
$M(n)$, and therefore defines a section $\alpha(m)$ of
$\mathcal{A}(M(n))=\mathcal{A}(M)(n)$ (compare
\hyperref[fac:III:2:no:58:prop:4]{Proposition~4}). This defines
$\alpha:M\to\Gamma(\mathcal{A}(M))$, and it is immediate that $\alpha$
is a homomorphism.

\phantomsection\label{fac:III:2:no:59:def:beta}
\textsc{Definition of $\beta$.}
Let $\mathcal{F}$ be an algebraic sheaf on $X$, and let $s/Q$ be an
element of $\Gamma(\mathcal{F})_x$, with
$s\in\Gamma(X,\mathcal{F}(n))$, $Q\in S_n$, and $Q(x)\ne0$. The
function $1/Q$ is homogeneous of degree $-n$ and regular at $x$, hence
is a section of $\mathcal{O}(-n)$ in a neighborhood of $x$. It follows
that $1/Q\otimes s$ is a section of
$\mathcal{O}(-n)\otimes\mathcal{F}(n)=\mathcal{F}$ in a neighborhood
of $x$, and thus defines an element of $\mathcal{F}_x$. We denote this
element by $\beta_x(s/Q)$, since it depends only on $s/Q$. The map
$\beta_x$ can also be defined using the components $s_i$ of $s$: if
$x\in U_i$, then
\phantomsection\label{fac:III:2:no:59:formula:beta-x}
\[
\beta_x(s/Q)=(t_i^n/Q).s_i(x).
\]
The collection of homomorphisms $\beta_x$ defines the homomorphism
$\beta:\mathcal{A}(\Gamma(\mathcal{F}))\to\mathcal{F}$.

The homomorphisms $\alpha$ and $\beta$ are related by the following
propositions, which follow by a direct calculation.

\phantomsection\label{fac:III:2:no:59:prop:7}
\textsc{Proposition 7.} \textit{Let $M$ be a graded $S$-module. The
composite of the homomorphisms}
\phantomsection\label{fac:III:2:no:59:formula:prop7-composite}
\[
\mathcal{A}(M)\longrightarrow
\mathcal{A}(\Gamma(\mathcal{A}(M)))\longrightarrow\mathcal{A}(M)
\]
\textit{is the identity.}

(The first homomorphism is defined by
$\alpha:M\to\Gamma(\mathcal{A}(M))$, and the second is $\beta$, applied
to $\mathcal{F}=\mathcal{A}(M)$.)

\phantomsection\label{fac:III:2:no:59:prop:8}
\textsc{Proposition 8.} \textit{Let $\mathcal{F}$ be an algebraic sheaf
on $X$. The composite of the homomorphisms}
\phantomsection\label{fac:III:2:no:59:formula:prop8-composite}
\[
\Gamma(\mathcal{F})\longrightarrow
\Gamma(\mathcal{A}(\Gamma(\mathcal{F})))
\longrightarrow\Gamma(\mathcal{F})
\]
\textit{is the identity.}

(The first homomorphism is $\alpha$, applied to
$M=\Gamma(\mathcal{F})$, while the second is defined by
$\beta:\mathcal{A}(\Gamma(\mathcal{F}))\to\mathcal{F}$.)

\phantomsection\label{fac:III:2:no:59:forecast:no:65}
We shall show in \hyperref[fac:III:3:no:65]{no.~65} that
$\beta:\mathcal{A}(\Gamma(\mathcal{F}))\to\mathcal{F}$ is bijective if
$\mathcal{F}$ is coherent, and that
$\alpha:M\to\Gamma(\mathcal{A}(M))$ is
\hyperref[fac:III:2:no:56:def:C-morphism]{$\mathcal{C}$-bijective} if
$M$ \hyperref[fac:III:2:no:56:def:TF]{satisfies condition (TF)}.
